

El análisis experimental confirmó que las cotas teóricas deben contrastarse con los factores constantes de implementación en escenarios reales. En la multiplicación de matrices, la sobrecarga recursiva y el elevado consumo de memoria de Strassen ($4.270$ KB frente a $232$ KB en $N=1024$) retrasan su ventaja asintótica, resultando el algoritmo Naive sistemáticamente más rápido en todas las dimensiones evaluadas. Para los algoritmos de ordenamiento, \texttt{std::sort} y Merge Sort mostraron una estricta estabilidad en $O(N \log N)$ bajo cualquier configuración, mientras que Quick Sort evidenció su sensibilidad al dominio numérico, degradándose drásticamente en $D_1$ hasta rozar los $300.000$ ms debido a la presencia masiva de duplicados. Por último, Patience Sort exhibió una fuerte dependencia de la distribución inicial, registrando su peor desempeño en arreglos ordenados ascendentemente al forzar la instanciación de múltiples pilas individuales.
